\section{Hermann Keyserling}

\begin{quotationx}
This is the part on \textbf{Julius Evola}`s interpretation of \textbf{Hermann Keyserling}, originally published in \emph{Saggi sull'Idealismo Magico}. The School of Wisdom\footnote{\url{https://www.schoolofwisdom.com/}} referred to was established by Keyserling in 1920 and is apparently still active today.

Keyserling bases his philosophy on ``understanding'', which is ``metarational''. Specifically, it is not something that can be given to someone nor compelled by a logical process. It should be familiar to all readers. For example, you may have had an ``aha'' experience, where you suddenly ``get it''. Or it could be as simple as ``getting a joke'' ... you all know someone who failed to comprehend the punchline! If you consider such moments, you will see how not understanding changed to understanding. This is a free spontaneous act, since nothing caused it, not even neurochemical reactions.

The second point is that meaning is always in reference to a Self, not in the object. The exterior world takes on quite a different feeling once it is realized the extent that the Self understands it and invests it with meaning. Matter, then, is the privation of meaning, since it belongs to necessity and understanding is free.

In the second part, we will see that through understanding and meaning, the world of representations is created. This is not to say that the material world is such, insofar as it is privation, but, rather, there is a large part of social life that is not matter, but rather the creation of the mind. This is the will to power, and it depends on reaching increasingly deeper levels of understanding.

This is not much different from what Boris Mouravieff says: ``Esoteric Tradition teaches that any civilization is none other than a projection of the consciousness of the I's of elite man onto the exterior world.'' Now if the elite is made up of Sages, as Keyserling describes, the projection will not be that of the chaotic I's of the personality, but rather of the Real I.

\end{quotationx}


Hermann Keyserling created tendencies that gave rise to the so-called ``School of Wisdom'' in Germany and are interesting for this reason: they aim to transfigure the intellectual synthesis itself in the metarational principle of freedom and therefore to shift the center of the I from the \emph{vis a tergo} of the rational imperative and every necessity in general, without suppressing the plane of concrete experience.

The key to Keyserling's views is the phenomenon of understanding. He contemplates the point in which the I says to himself: ``I understood.'' It is essentially a point of spontaneity, freedom, and interiority: there is no way of compelling one to understanding. It on the other hand has a mystico-illuminative character: it seeks to \emph{live} the moment of ``meaning'', not of this or that meaning, but rather of meaning in general, of the pure element of the understanding conditioning every understanding, then one will feel that something ineffable shines in it, something that, even if contending with it and being supported on it, absolutely transcends the whole of the means and forms from which it was propitiated. This mystical moment of pure understanding is the moment of the spirit. Certainly, a meaning exists only in connection to a certain form or nature that expresses it --- but the inverse is also true, and that is that being understood, interiorized, is always the condition so that anything has existence for the I and that therefore the synthetic function of understanding is, directly, the absolute \emph{prius}, the base, or \emph{apriori} of each one of our experiences, if only in accordance with somewhat diverse levels.

Since it is absurd to speak of a meaning existing in things independently of the I (nothing has meaning through itself but everything can acquire it) and since understanding is always and essentially interwoven with spontaneity and interiority, it appears clear that by that, Keyserling defines a function, capable of reaffirming, on the entire ambit of human experience, the principle of the free and creator I. Since --- and this is Keyserling's improvement, previously sketched out by Novalis, on Kant's synthetic a priori which, as everyone knows, is preconscious, impersonal, and abstractly intellectual --- there understanding is the same as an unconditioned and immanent power of the real I, it is not a concept, but an effective element of interiority. It follows that the entire world assumes the character of an expressive means, of a symbolic material that the I must invest, animate, and almost recreate with the act of his understanding.

The distinction between nature and spirit is thus reduced to that between a half expressive abstraction and half expressive entirely resolved in the actuality of meaning. Matter or necessity would only be the \emph{privation} of meaning, the mindless opaque ``letter'' to oneself. Nevertheless, even reduced to this form, such a distinction must be explained. It is that meaning, in its essence of deep subjectivity and in law, is absolute freedom: but, in expressing itself, it cannot be crystalized into a given nonconvertible body, it cannot make that unconditioned principle fall into necessity (that is, in a mechanical discourse) in which it comes to fruition --- as in a type of enlightenment --- in the pure moment of self-determination, of the creative conception balanced between the ``not yet'' of the possible and the form in which the possible asserts itself. The half expressive abstraction, or nature and necessity, is constituted by nothing other than the preceding processes of self-expression already exhausted and sustained through mechanical repetition; from which however, the spirit rises in self-conception on the basis of a deeper meaning, which it then expresses and incarnates utilizing precisely that matter in which it is coagulated and which rendered its preceding freedom mechanical --- and the process is continued toward an always deeper meaning or interiority, correlative to an always richer, organic, and articulate body of expression. From this rises the concept that the various natural or historical laws are simple laws of grammar and syntax, which those who live in the deep level of meaning do not need to deny, but only to dominate interiorly, just as the artist dominates the material in which he incorporates his creation.

That brings up the problem of the human ``type''. It, according to Keyserling, is not to be pushed back onto art, religion, or philosophy. Artists are, typically, mediums: that greatness that speaks in, or by means of, them almost never coincides with their personal knowledge. As for religion, it must be excluded because, as such, it introduces a principle of authority and dogmatism on the one hand, of dependence and passivity on the other, that is incompatible with the character of autonomous, individual assertion, that was connected to ``meaning''. Still less can it be pushed back to the thinker, who is held fixed in a world of abstract concepts, alien to reality and separated from the depths of the creator I. It is better instead to connect it to the ancient concept of \emph{Wisdom}, when it is understood as a synthesis of life and science in the individual unity of the creator. What is essential, is that the I does not become a slave of abstract knowledge, but produces it interiorly in living reality; that he survives in an ideality no more than an abyss separated from concreate reality, but rather he places himself outside the world and \emph{expresses} himself in it, incarnates fully in it to the extent he conceives the deep level of meaning. Hence, the type of the Sage becomes moreover that of the Lord: and that clearly not in regard to naked power, but rather in regard to the principle that is interiorly superior to the totality of life, that commands it and molds from above the power of freedom. And how much deeper is the plane in which understanding is realized, so much more perfectly and fully the I dominates and commands the whole of the various forces (natural, social, etc.) that it reclaims in that way as the matter of his ``language''.

From this development it becomes clear that Keyserling speaks of ``meaning'' and ``understanding'' in a rather metaphoric way, in order to give a suggestion of a certain function, which he then, to tell the truth, reaffirms in the totality of all concrete powers of internal and external experience.

\begin{quotationx}
This is the second and final installment of \textbf{Julius Evola}`s commentary on \textbf{Hermann Keyserling} from \emph{Saggi sull'Idealismo Magico}. Evola refers to Keyserling's ``brilliant interpretation of the function of meaning, according to which understanding is removed from the rational and peripheral plane and compenetrated with the principle of deep self-realization and power.''

Some of the main points to note:

\begin{itemize}


\item There are deeper levels of consciousness from which the Self has more power over the physical world.

\item The physical world is not self-sufficient, but is related to a spiritual power for its ultimate explanation

\item The accumulation of facts and knowledge is not the important thing. Rather, it is the deepening of one's level of consciousness

\item Meaning is really the self-possession of the person

\item The physical world is not destiny, it does not consist of brute facts

\item Reality is in fact malleable, perfectly reflecting what the Self conceives as its meaning

\item The individual must accept full responsibility for his world

\item In that way, he overcomes fate and becomes free
\end{itemize}
\end{quotationx}


Clearly, the problem of knowledge in the doctrine of ``meaning'' is strictly connected to the problem of power, which Keyserling resolves by means of the theory of the levels of consciousness. As we previously noted as characteristic of idealism, the premise is that every objective thing depends on a subjective thing, that things are what we are or, better, what we posit to ourselves. To will to dominate the external world, i.e., acting in the level of the mass and physical determinisms, is an impossible and contradictory assumption: but things are quite different when the I brings its own action into the deep level of transcendent causes, into the sphere of ``meaning'' that is that incorporeity that conditions the corporeal and that is conditioned by nothing. Human liberty in its highest aspect consists in the idea that where the focus of our consciousness falls depends on us. We can posit this focus in the world of the phenomenon where there is no place for a real initiative, or else at the point of the original creative function, whose principle is freedom and \emph{possibility}. In other words: every phenomenon never constitutes a last instance, but presupposes a spiritual power to which its explanation and consistency refers back: when we posit ourselves with this power in relation to the other, when we do not \emph{understand} it, it appears to us as unbending fate; when instead it is reaffirmed to us in that divine spark that remains in the individual and that is the deep source of his life, that world that first held us in iron slavery becomes our instrument for a silent transformation.

Then it will be able to determine unconditionally the form in which reality will have to appear to us, to actually live the given not as a last instance, but as one that is malleable and reflects passively what the I conceives in the region of meaning. ``The representation creates reality, and not vice versa''; ``the faculty of representation is unconditioned''; the I can, by means of a movement of the level of consciousness, possessing itself in this faculty''---such are the principles of Keyserling's doctrine of power similar to those characteristic of esoterism and certain western schools of magic.

The fundamental point is added to that, even if not so distinctly affirmed like Michelstaedter: namely, the individual must rise to the awareness of absolute responsibility, he must make himself adequate to his own life and not only in the order of the subjective properly called, but also in the order of the cosmic and the universal. The I must make himself the last cause and must be able to take the weight of world responsibility on himself, without trying to push it back onto another; since only on condition of assuming it, can the person hope to go beyond destiny to freedom.

The principle merit of the ``School of Wisdom'' is that it affirmed that what is important is not to procure new knowledge or experiences, but rather to change the \emph{level} of one's own personality, to bring the center of the I onto another plane or dimension in which the relationship to one's activity is that of possession and unconditioned determination. In that he agrees fully with magical idealism, and still more so when he suggests such a plane as that of the Lord. The concept of ``meaning'' instead demands a greater individuation. In fact the question must be asked: what is the \emph{meaning} of this ``meaning''? Since it certainly has something to do with its general mystical and self-creative aspect: but when the dialectic of expression mentioned above is grafted onto it, things become complicated. ``Meaning'' understood as a fate of self-expression, objectifying itself, reasserting itself in new expressions, etc., is what can be less comprehensible and in no way surrenders to the brute given of dialectic laws to which rationalism is subjected to.

Even here, freedom must certainly be asserted beyond the entire process. On the other hand, in order to organically connect the same process---understood as the ever deeper self-penetration of meaning---to the principle of the Lord, it is necessary to understand this dynamically, i.e., to place it properly in the \emph{interval} in which the I transcends his freedom which became an objectivity, in order to reassert a principle eternally irreducible to \emph{being}. Then ``understanding'', ``meaning'' should have self-possession for its deepest meaning, the self-realizing of the I in pure actual essence in an ever more perfect way. \emph{Autarchy} should thus be the key to the process, the quality in which original freedom willed itself---consequently, Keyserling clouds over when he points to the principle that is necessary to make oneself the last reason, that it is necessary to rise to the meaning of world responsibility, to take destiny onto oneself, to make oneself adequate to it and to resolve it in the principle of the Lord. Besides, once that determinism is established, the rigidity is what constitutes the antithetical moment since as matter for the expression, the reassertion of the principle of meaning---which is freedom---follows from it or it can have only a negative character---of \emph{derealization}, of agitation, of dissolution of every necessity in the contingent and in the easily moveable---so that the idea of a substrate of material necessity, of laws, whether restricted only to the only ``grammar'' or ``language'' (the same languages moreover change, transforming themselves according to various needs), cannot in the final instance be maintained, and the idea of a world system entirely compenetrated by the function of the ``meaning'' is enriched with the idea of the ``body of freedom'' that is also the ``body of negation''. With Heraclitus, Novalis, Bhagavan Das, and the conception here sustained, it is necessary to understand the world process as a \emph{burning} and its immanent goal as the accomplishment of the I as \emph{entity of pure negativity}---e.g., of pure contingency and absolute possession.

If we wanted, we could still make a few and more serious objections to Keyserling, to the asystematic and ``symphonic'' character of his expositions: in the actual form specified, in regards to history, by which he does not mean completely \emph{ideality}, and that therefore he often treats from an empirical point of view, really transcendent to the plane of ``meaning''. So that while on one hand he affirms that the historic fact does not create ``meaning'' but that ``meaning'' creates the historic fact, on the other, like Benedetto Croce, he in fact nevertheless recognizes given historical determinisms that impinge, epoch by epoch, the uneliminable conditions so that meaning can incarnate itself and make itself concrete. We may think of Keyserling's principle that meaning exists only insofar as it is expressed, and one will notice here how dangerous the position is and how instead---since he holds firm to the principle, that nothing is that is understood in a certain measure and that understanding has freedom for an inseparable attribute---the right way would be to deny the \emph{real} anteriority of any historical condition whatsoever, and therefore of the entire past or antecedent history to make absolutely unformed matter, whose identification is not \emph{given} or \emph{found}, but is pinned only to the metahistoric unconditioned self-determination of meaning---i.e., to make something conditioned from it and no longer something conditioning.

But this is not the place to repeat myself. It was important only to point out Keyserling's brilliant interpretation of the function of ``meaning'', according to which understanding is removed from the rational and peripheral plane and compenetrated with the principle of deep self-realization and power.


\flrightit{Posted on 2014-07-29 by Cologero}

